\chapter[Branch-line reductions and primitive modular characteristics]{Branch-line reductions, primitive modular characteristics, and the intrinsic invariant ladder}
\label{app:branch-line-reductions}

The full modular Maurer--Cartan element $\Theta_\cA$ lives in an
infinite-dimensional convolution algebra.  A \emph{branch-line
reduction} restricts~$\Theta_\cA$ to a configuration of minimal
complexity (two-point branch lines, rank-one spectral sectors,
genus-$2$ transparency windows) and proves structural theorems
inside the resulting finite-dimensional quotient.

The \emph{primitive modular characteristics} are the irreducible
components of~$\Theta_\cA$ under the clutching maps of
$\overline{\mathcal{M}}_{g,n}$: they are to the shadow obstruction
tower what primitive elements are to a Hopf algebra.  The intrinsic
invariant ladder
\[
\kappa(\cA)
\rightsquigarrow
\Delta_{\cA}(x)
\rightsquigarrow
\Theta_{\cA}.
\]
The load-bearing theorems are the scalar modular characteristic
(Theorem~\ref{thm:modular-characteristic}), the spectral characteristic
(Theorem~\ref{thm:spectral-characteristic}), and the universal modular
Maurer--Cartan theorem
(Theorems~\ref{thm:universal-theta} and~\ref{thm:mc2-full-resolution}).
The purpose of the appendix is not to replace the full MC2 package by a
smaller surrogate.  It is to extract exact finite-dimensional quotients
of that package and prove structural theorems inside those quotients.
Those quotient theorems clarify which parts of the higher-genus story are
already forced by the scalar and spectral shadows, and which parts still
belong to the genuinely infinite cyclic \(L_\infty\) geometry.

\medskip
\noindent\emph{Two corrections organize the whole discussion.}
First, the natural intrinsic ladder is
\(\kappa\), \(\Delta\), \(\Theta\), not a comparison theorem such as
Drinfeld--Kohno.  The spectral discriminant is the first non-scalar
invariant built purely from the modular Koszul package, whereas DK-type
comparison theorems sit later in the theorem graph.
Second, the universal class \(\Theta_{\cA}\) is available by
Theorems~\ref{thm:universal-theta} and~\ref{thm:mc2-full-resolution}.
For one-channel families, the minimal class lies on a single line
\(\eta\otimes\Gamma_{\cA}\); on the proved scalar lane this further
specializes to \(\kappa(\cA)\eta\otimes\Lambda\).
The non-scalar content is organized by the shadow
Postnikov tower (Definition~\ref{def:shadow-postnikov-tower}),
with finite-order shadows proved through arity~$4$.
The new constructions below therefore have a precise role: they are
shadow and quotient theorems \emph{inside} the proved
finite-order modular package.

\medskip
\noindent\emph{What is new here.}
The appendix proves five families of results.
\begin{enumerate}
\item It constructs the scalar Maurer--Cartan skeleton and the spectral
  cumulant hierarchy as canonical logarithmic shadows of
  \(\Theta_{\cA}\).
\item It proves an abstract filtered obstruction theorem showing that
  the first nontrivial correction to a scalar shadow is necessarily
  traceless and quadratic.
\item It proves that the positive-weight Virasoro sector is contractible
  in the effective \(L_0\)-model, so the first non-scalar content is
  forced into finite reduced spectral sectors.
\item It proves a genus-\(2\) transparency theorem for every rank-two
  branch-line quotient with scalar square.
\item It computes the first primitive traceless coefficients in the pure
  branch and top-Hodge primitive quotients, including the universal
  genus-\(3\) constant~\(7\) on the shared
  \(\widehat{\mathfrak{sl}}_2\)/Virasoro/\(\beta\gamma\) spectral
  sheet.
\end{enumerate}

\begin{construction}[Branch lines as shadow depth loci]
\label{constr:branch-line-shadow-depth}
Branch-line reductions are the geometric loci where the shadow
depth $r_{\max}(\cA)$ transitions between archetype classes.
The weight-changing line in the $\beta\gamma$ system
(class~$\mathsf{C}$, $r_{\max} = 4$) degenerates to the
Heisenberg line (class~$\mathsf{G}$, $r_{\max} = 2$) when
the quartic contact invariant $\mu_{\beta\gamma}$ vanishes
(Corollary~\ref{cor:nms-betagamma-mu-vanishing}).
The general mechanism has three layers.
\begin{enumerate}[label=\textup{(\roman*)}]
\item The shadow depth $r_{\max}(\cA)$ is the smallest
  integer~$r$ such that the obstruction class
  $o_{r+1}(\cA) \in H^2(F^{r+1}\mathfrak{g}/F^{r+2}\mathfrak{g}, d_2)$
  vanishes identically, or $\infty$ if no such~$r$ exists
  (Definition~\ref{def:shadow-postnikov-tower}).
  It classifies chiral algebras into four archetype
  classes: $\mathsf{G}$ ($r_{\max} = 2$, Gaussian),
  $\mathsf{L}$ ($r_{\max} = 3$, Lie/tree),
  $\mathsf{C}$ ($r_{\max} = 4$, contact/quartic),
  $\mathsf{M}$ ($r_{\max} = \infty$, mixed).
\item A branch line is a one-parameter family in which the
  shadow depth drops: there exists a critical value of the
  spectral parameter at which $o_{r_{\max}}(\cA)$
  becomes cohomologous to zero, reducing the effective
  shadow depth by at least one.  The primitive modular
  characteristics are the irreducible components of
  $\Theta_\cA$ that survive on the branch line; those
  that vanish are precisely the obstructions that the
  branch-line reduction removes.
\item For the $\beta\gamma$ system, the
  weight-changing branch line is the locus where the
  quartic contact invariant
  $\mu_{\beta\gamma}
  = \langle \eta, m_3(\eta,\eta,\eta)\rangle$
  vanishes.  At this locus the shadow obstruction tower truncates from
  arity~$4$ to arity~$2$, and the residual modular
  characteristic is the scalar~$\kappa(\beta\gamma)
  = 1$, matching the Heisenberg value at rank~$1$.
  This is the content of the rank-one abelian rigidity
  argument of Corollary~\ref{cor:nms-betagamma-mu-vanishing}.
\end{enumerate}
\end{construction}

\section{The invariant ladder from first principles}
\label{sec:blr-invariant-ladder}

The modular characteristic hierarchy has three logarithmic levels and
three different kinds of geometry.
\begin{enumerate}[label=\textup{(\roman*)}]
\item The scalar level is governed by the single number
  \(\kappa(\cA)\) and the tautological line
  \(\kappa(\cA)\lambda_g\) on \(\overline{\mathcal M}_g\)
  (Theorem~\ref{thm:modular-characteristic}).
\item The spectral level is governed by the spectral discriminant
  \(\Delta_{\cA}(x)\), equivalently by the branch transport operator in
  the finite-rank spectral model
  (Theorem~\ref{thm:spectral-characteristic}).
\item The full level is governed by the cyclic deformation complex
  \(\Defcyc(\cA)\) and the universal Maurer--Cartan class
  \(\Theta_{\cA}\)
  (Theorems~\ref{thm:universal-theta}
  and~\ref{thm:mc2-full-resolution}).
\end{enumerate}
The point of view of this appendix is that the scalar and spectral levels
are not secondary bookkeeping devices.  They are exact quotient faces of
\(\Theta_{\cA}\), and they can therefore be studied as mathematical
objects in their own right.

\begin{definition}[Scalar Maurer--Cartan skeleton]
\label{def:scalar-mc-skeleton}
Let \(\cA\) be a modular Koszul chiral algebra.  Define the abelian dg
Lie algebra
\[
\mathfrak{s}_{\cA}:=\prod_{g\geq 1}\mathbf Q\,u_g,
\qquad
|u_g|=1-2g,
\qquad
d=0,
\qquad
[-,-]=0.
\]
The \emph{scalar Maurer--Cartan skeleton} of \(\cA\) is the element
\begin{equation}
\label{eq:scalar-mc-skeleton}
\Theta^{\mathrm{sc}}_{\cA}
:=
\sum_{g\geq 1}
\kappa(\cA)
\,u_g\otimes \lambda_g
\in
\mathfrak{s}_{\cA}
\widehat\otimes
\prod_{g\geq 1}H^{2g}(\overline{\mathcal M}_g,\mathbf Q).
\end{equation}
\end{definition}

\begin{proposition}[The scalar shadow is an abelian MC element; \ClaimStatusProvedHere]
\label{prop:scalar-mc-skeleton}
For every modular Koszul chiral algebra~\(\cA\), the class
\(\Theta^{\mathrm{sc}}_{\cA}\) of
\eqref{eq:scalar-mc-skeleton} is a Maurer--Cartan element of the
abelian dg Lie algebra \(\mathfrak{s}_{\cA}\).  Moreover, it is exactly
 the image of \(\Theta_{\cA}\) under the scalar trace map of
Theorem~\ref{thm:universal-theta}.
\end{proposition}

\begin{proof}
Each summand
\(u_g\otimes \lambda_g\) has total degree
\((1-2g)+2g=1\), so
\(\Theta^{\mathrm{sc}}_{\cA}\) has Maurer--Cartan degree~\(1\).
Since \(\mathfrak{s}_{\cA}\) is abelian and has zero differential, the
MC equation reduces to closedness, which is immediate because the
classes \(\lambda_g\) are closed.  The trace identity is precisely
clause~\textup{(i)} of Theorem~\ref{thm:universal-theta}.
\end{proof}

\begin{remark}[The scalar ladder is a quotient, not a replacement]
\label{rem:scalar-skeleton-not-replacement}
The existence of \(\Theta^{\mathrm{sc}}_{\cA}\) does not compete with the
full theorem.  It identifies the maximal abelian quotient of the full
modular deformation problem.  Every finite or filtered approximation to
\(\Theta_{\cA}\) that respects the scalar trace must reduce to
\(\Theta^{\mathrm{sc}}_{\cA}\) after projecting away all brackets and all
non-scalar channels.
\end{remark}

\section{Spectral cumulants and the connected potential}
\label{sec:blr-spectral-cumulants}

Theorem~\ref{thm:spectral-characteristic} isolates the first intrinsic
non-scalar invariant: the spectral discriminant.  The next step is to
extract from it a canonical infinite hierarchy of characteristic
numbers.

\begin{definition}[Spectral cumulants and connected potential]
\label{def:spectral-cumulants}
Let \(\cA\) be a modular Koszul chiral algebra for which the spectral
characteristic is realized by a finite-rank branch transport operator
\(T_{\mathrm{br}}\) with
\[
\Delta_{\cA}(x)=\det(1-xT_{\mathrm{br}}).
\]
Define the \emph{spectral cumulants}
\[
\mathfrak c_m(\cA):=\operatorname{Tr}(T_{\mathrm{br}}^m),
\qquad m\geq 1,
\]
and the \emph{connected spectral potential}
\begin{equation}
\label{eq:connected-spectral-potential}
\Psi_{\cA}(x):=-\log \Delta_{\cA}(x).
\end{equation}
\end{definition}

\begin{theorem}[Spectral cumulant hierarchy; \ClaimStatusProvedHere]
\label{thm:spectral-cumulant-hierarchy}
The connected potential satisfies
\begin{equation}
\label{eq:spectral-cumulant-expansion}
\Psi_{\cA}(x)
=
\sum_{m\geq 1}
\frac{\mathfrak c_m(\cA)}{m}
\,x^m.
\end{equation}
Consequently:
\begin{enumerate}[label=\textup{(\roman*)}]
\item each \(\mathfrak c_m(\cA)\) is a canonical spectral invariant;
\item the discriminant \(\Delta_{\cA}(x)\) is determined by finitely
  many cumulants once its degree is fixed;
\item in the shared rank-two spectral sheet
  \(\Delta(x)=(1-3x)(1+x)\) of
  Theorem~\ref{thm:ds-bar-gf-discriminant}, one has
  \begin{equation}
  \label{eq:shared-sheet-cumulants}
  \mathfrak c_m=3^m+(-1)^m.
  \end{equation}
\end{enumerate}
\end{theorem}

\begin{proof}
Using \(\Delta_{\cA}(x)=\det(1-xT_{\mathrm{br}})\), one computes
\[
-\log\Delta_{\cA}(x)
=
-\log\det(1-xT_{\mathrm{br}})
=
-\operatorname{Tr}\log(1-xT_{\mathrm{br}})
=
\sum_{m\geq 1}
\frac{\operatorname{Tr}(T_{\mathrm{br}}^m)}{m}
\,x^m,
\]
which is \eqref{eq:spectral-cumulant-expansion}.  This proves
\textup{(i)}.  For \textup{(ii)}, write
\(\Delta_{\cA}(x)=1-e_1x+e_2x^2-\cdots+(-1)^d e_dx^d\).  The Newton
identities recover the elementary symmetric polynomials \(e_i\) from the
power sums \(\mathfrak c_1,\dots,\mathfrak c_d\), so those first \(d\)
cumulants determine \(\Delta_{\cA}\).  Finally, in the shared sheet,
\(T_{\mathrm{br}}\) has eigenvalues \(3\) and \(-1\), hence
\eqref{eq:shared-sheet-cumulants}.
\end{proof}

\begin{remark}[Connected cycle interpretation]
\label{rem:connected-cycle-interpretation}
Formula~\eqref{eq:spectral-cumulant-expansion} is the connected-cycle
expansion of the spectral transport.  The coefficient
\(\mathfrak c_m(\cA)/m\) is the cyclic trace of length-\(m\) branch
transport, divided by the number of cyclic choices of basepoint.  In
particular, \(\Psi_{\cA}(x)\) is the connected one-loop spectral shadow
of the full modular Maurer--Cartan class.
\end{remark}

\section{An abstract filtered obstruction theorem}
\label{sec:blr-filtered-obstruction}

The full MC2 theorem resolves the existence of \(\Theta_{\cA}\).  The
next theorem clarifies a different question: where can a first
nontrivial correction appear in a trace-compatible quotient or
truncation of a Maurer--Cartan problem?

\begin{definition}[Trace-compatible filtered dg Lie algebra]
\label{def:trace-compatible-filtered-dglie}
A \emph{trace-compatible filtered dg Lie algebra} consists of:
\begin{enumerate}[label=\textup{(\roman*)}]
\item a complete descending filtration
  \(L=F^1L\supseteq F^2L\supseteq \cdots\)
  by dg Lie ideals;
\item an abelian dg Lie algebra \(S\);
\item a morphism of filtered dg Lie algebras
  \(\operatorname{tr}\colon L\to S\)
  such that \(\operatorname{tr}[x,y]=0\) for all \(x,y\in L\).
\end{enumerate}
Write \(L^\circ:=\ker(\operatorname{tr})\) for the traceless sector.
\end{definition}

\begin{theorem}[First obstruction is traceless and quadratic; \ClaimStatusProvedHere]
\label{thm:first-obstruction-traceless-quadratic}
Let \((L,F^\bullet,\operatorname{tr})\) be a trace-compatible filtered dg
Lie algebra, and let \(\theta_{\mathrm{sc}}\in S^1\) be a closed degree-\(1\)
element.  Suppose that \(\vartheta^{(1)}\in F^1L/F^2L\) is a first-order
lift of \(\theta_{\mathrm{sc}}\).  For any representative
\(\widetilde\vartheta_1\in F^1L\), define the curvature modulo \(F^3L\) by
\begin{equation}
\label{eq:first-curvature-mod-f3}
\Omega_2(\widetilde\vartheta_1)
:=
 d\widetilde\vartheta_1
 +\frac12[\widetilde\vartheta_1,\widetilde\vartheta_1]
 \mod F^3L.
\end{equation}
Then:
\begin{enumerate}[label=\textup{(\roman*)}]
\item \(\Omega_2(\widetilde\vartheta_1)\) defines a class
  \begin{equation}
  \label{eq:first-obstruction-class}
  o_2\in H^2(F^2L/F^3L)
  \end{equation}
  independent of the representative;
\item \(o_2=0\) if and only if \(\vartheta^{(1)}\) lifts to a
  Maurer--Cartan solution modulo \(F^3L\);
\item the class \(o_2\) lies in the traceless sector
  \(H^2(F^2L^\circ/F^3L^\circ)\).
\end{enumerate}
\end{theorem}

\begin{proof}
Because \(\vartheta^{(1)}\) is closed modulo \(F^2L\), the curvature of
\(\widetilde\vartheta_1\) lies in \(F^2L\).  A standard graded Jacobi
identity computation gives
\[
 d\Omega_2(\widetilde\vartheta_1)
 \equiv 0
 \pmod{F^3L},
\]
so \(\Omega_2\) is a cocycle modulo \(F^3L\).  If one changes the
representative by an element \(\alpha\in F^2L\), then
\[
\Omega_2(\widetilde\vartheta_1+\alpha)
-
\Omega_2(\widetilde\vartheta_1)
\equiv d\alpha
\pmod{F^3L},
\]
because every bracket involving \(\alpha\) lands in \(F^3L\).  This
proves \textup{(i)}.

For \textup{(ii)}, a correction
\(\vartheta_2\in F^2L/F^3L\) produces a Maurer--Cartan solution modulo
\(F^3L\) exactly when
\[
 d\vartheta_2 = -\Omega_2(\widetilde\vartheta_1)
 \quad \text{in }F^2L/F^3L.
\]
So the obstruction vanishes precisely when the class \(o_2\) vanishes.

Finally,
\[
\operatorname{tr}(\Omega_2(\widetilde\vartheta_1))
=
d\operatorname{tr}(\widetilde\vartheta_1)
+\frac12\operatorname{tr}[\widetilde\vartheta_1,\widetilde\vartheta_1]
=
d\theta_{\mathrm{sc}}+0=0.
\]
Hence the obstruction class lies in the kernel of the trace map.  This
is \textup{(iii)}.
\end{proof}

\begin{corollary}[Vanishing criterion for filtered lifts; \ClaimStatusProvedHere]
\label{cor:filtered-lift-vanishing}
In the situation of
Theorem~\ref{thm:first-obstruction-traceless-quadratic}, assume in
addition that
\[
H^2\bigl(F^nL^\circ/F^{n+1}L^\circ\bigr)=0
\qquad\text{for all }n\geq 2.
\]
Then every scalar Maurer--Cartan skeleton admits a lift to an honest
Maurer--Cartan element of~\(L\).  If moreover
\[
H^1\bigl(F^nL^\circ/F^{n+1}L^\circ\bigr)=0
\qquad\text{for all }n\geq 2,
\]
then the lift is unique up to gauge equivalence among all lifts with the
same scalar trace.
\end{corollary}

\begin{proof}
The proof is the standard inductive extension argument.  At stage
\(n\), the failure of a lift modulo \(F^{n+2}L\) defines a class in
\(H^2(F^{n+1}L/F^{n+2}L)\).  Repeating the trace computation from the
proof of
Theorem~\ref{thm:first-obstruction-traceless-quadratic}, one sees that
this class lies in the traceless sector.  The vanishing of
\(H^2\) kills the obstruction and produces a lift at the next stage.
Completeness of the filtration yields a limit Maurer--Cartan element.
The vanishing of \(H^1\) gives uniqueness by the usual gauge-induction
argument.
\end{proof}

\begin{remark}[Use after MC2]
\label{rem:filtered-lift-use-after-mc2}
In the present repository the full lift already exists, so
Corollary~\ref{cor:filtered-lift-vanishing} is not needed to prove
\(\Theta_{\cA}\).  Its role is diagnostic: any finite or spectral
approximation to \(\Theta_{\cA}\) begins to fail, if it fails at all,
in the traceless quadratic sector.
\end{remark}

\section{Virasoro positive-weight acyclicity and reduced spectral localization}
\label{sec:blr-vir-positive-weight}

This section isolates the first genuinely finite obstruction sector in
the Virasoro family.  The mechanism is the same one that already drives
the PBW calculations: positive conformal weight is invertible under the
\(L_0\)-action.

\begin{definition}[Effective positive-weight functor]
\label{def:effective-positive-weight-functor}
Let \(\bar V_{\mathrm{Vir}_c}=\bigoplus_{h\geq 2} M_h\) be the
augmentation ideal of the Virasoro vacuum algebra, graded by conformal
weight.  An \emph{effective positive-weight functor} on
\(\bar V_{\mathrm{Vir}_c}\) is a finite polynomial Schur functor built
from tensors, symmetric powers, exterior powers, and finite-dimensional
auxiliary spaces, with the condition that every homogeneous monomial has
strictly positive total conformal weight.

If \(\mathcal F(\bar V_{\mathrm{Vir}_c})\) is such a functor, define the
\emph{effective differential}
\begin{equation}
\label{eq:effective-L0-differential}
 d_{\mathrm{eff}}:=\mathcal F(L_0),
\end{equation}
where \(L_0\) acts diagonally by conformal weight.
\end{definition}

\begin{lemma}[Positive-weight contraction; \ClaimStatusProvedHere]
\label{lem:positive-weight-contraction}
Let \(W=\bigoplus_{m>0}W_m\) be a positively graded vector space, and let
\(D\) act on \(W_m\) by \(m\,\mathrm{id}\).  Then the complex
\((W,D)\) is contractible.
\end{lemma}

\begin{proof}
Define \(h\colon W\to W\) by
\(h|_{W_m}=m^{-1}\mathrm{id}\).
Then on \(W_m\) one has
\[
Dh+hD=\mathrm{id}_{W_m}.
\]
Hence \(\mathrm{id}_W\) is null-homotopic.
\end{proof}

\begin{theorem}[Positive-weight Virasoro sectors are acyclic; \ClaimStatusProvedHere]
\label{thm:vir-positive-weight-acyclic}
For every effective positive-weight functor
\(\mathcal F(\bar V_{\mathrm{Vir}_c})\), the complex
\((\mathcal F(\bar V_{\mathrm{Vir}_c}),d_{\mathrm{eff}})\) is
contractible.  In particular,
\[
H^\bullet\bigl(\mathcal F(\bar V_{\mathrm{Vir}_c}),d_{\mathrm{eff}}\bigr)=0.
\]
\end{theorem}

\begin{proof}
Every homogeneous monomial of
\(\mathcal F(\bar V_{\mathrm{Vir}_c})\) has total conformal weight
\(m>0\).  On the weight-\(m\) subspace,
\(d_{\mathrm{eff}}\) acts by \(m\,\mathrm{id}\).  Apply
Lemma~\ref{lem:positive-weight-contraction} weight by weight.
\end{proof}

\begin{corollary}[Localization to reduced spectral sectors; \ClaimStatusProvedHere]
\label{cor:vir-localization-reduced-spectral}
In every quotient of the Virasoro modular deformation problem whose
associated graded is built from effective positive-weight functors and a
finite-dimensional reduced spectral sector, the positive-weight part is
acyclic.  Consequently the first non-scalar quotient data are forced into
finite reduced spectral sectors.
\end{corollary}

\begin{proof}
The associated graded of the positive-weight part is acyclic by
Theorem~\ref{thm:vir-positive-weight-acyclic}.  Hence every spectral
sequence argument collapses there, leaving only the reduced finite-rank
sector.
\end{proof}

\begin{remark}[The reduced line on the shared sheet]
\label{rem:reduced-line-shared-sheet}
On the shared sheet
\(\Delta(x)=(1-3x)(1+x)\) of
Theorem~\ref{thm:ds-bar-gf-discriminant}, the reduced spectral algebra
is two-dimensional.  Its traceless part is therefore one-dimensional.
Any reduced traceless coefficient in a one-channel quotient must lie on a
single line.  The genus-\(2\) coefficient in the branch-line quotient is
therefore forced, a priori, to have the form
\(\mu_2(c)\,\tau\otimes \lambda_2\).
\end{remark}

\begin{proposition}[Odd-sheet rigidity for one-line reductions; \ClaimStatusProvedHere]
\label{prop:odd-sheet-rigidity}
Let \(t=k+2\) and write
\begin{equation}
\label{eq:vir-sheet-coordinate}
 c=1-\frac{6(k+1)^2}{k+2}=13-6\bigl(t+t^{-1}\bigr).
\end{equation}
Suppose that a reduced traceless Virasoro coefficient is represented on
the shared spectral line by \(m(t)\tau\), where:
\begin{enumerate}[label=\textup{(\roman*)}]
\item \(m(t)=m(t^{-1})\) \emph{(sheet descent)};
\item the same-family involution \(k\mapsto -k-4\), equivalently
  \(t\mapsto -t\), acts by \(\tau\mapsto -\tau\)
  \emph{(odd spectral normalization)}.
\end{enumerate}
Then
\begin{equation}
\label{eq:odd-about-13-general}
 m(t)=\bigl(t+t^{-1}\bigr)
 Q\!\left(\bigl(t+t^{-1}\bigr)^2\right)
\end{equation}
for a unique function \(Q\).  Equivalently, the descended coefficient is
of the form
\begin{equation}
\label{eq:mu-odd-about-13}
 \mu(c)=(c-13)P\bigl((c-13)^2\bigr)
\end{equation}
for a unique function \(P\).
\end{proposition}

\begin{proof}
Set \(u=t+t^{-1}\), so \(c=13-6u\).  The sheet-descent assumption says
that \(m\) depends only on \(u\).  Under the same-family involution,
\(t\mapsto -t\), hence \(u\mapsto -u\).  Because the spectral line is
odd, the transformed coefficient must satisfy
\(m(-t)=-m(t)\), so the descended function of \(u\) is odd.  Every odd
function of \(u\) has the form \(uQ(u^2)\), proving
\eqref{eq:odd-about-13-general}.  Substituting
\(u=(13-c)/6\) yields \eqref{eq:mu-odd-about-13}.
\end{proof}

\begin{corollary}[The genus-\(2\) one-line coefficient is centered at \texorpdfstring{$13$}{13}; \ClaimStatusProvedHere]
\label{cor:mu2-centered-at-13}
Every genus-\(2\) reduced traceless coefficient on the shared spectral
line that obeys the hypotheses of
Proposition~\ref{prop:odd-sheet-rigidity} satisfies
\[
\mu_2(26-c)=-\mu_2(c),
\qquad
\mu_2(13)=0.
\]
\end{corollary}

\section{Rank-two branch-line reductions}
\label{sec:blr-rank-two-branch}

The previous sections motivate a general finite-dimensional quotient.
The right level of generality is a rank-two spectral sheet.

\begin{construction}[Rank-two branch algebra]
\label{const:rank-two-branch-algebra}
Fix distinct scalars \(\lambda_+,\lambda_-\in \mathbf C\), and define
\begin{equation}
\label{eq:rank-two-branch-algebra}
 R_{\lambda_+,\lambda_-}
 :=
 \mathbf C[K]\big/\bigl((K-\lambda_+)(K-\lambda_-)\bigr).
\end{equation}
Set
\[
 c:=\frac{\lambda_++\lambda_-}{2},
 \qquad
 \delta:=\frac{\lambda_+-\lambda_-}{2},
 \qquad
 \tau:=K-c.
\]
Then
\begin{equation}
\label{eq:tau-square-general}
 R_{\lambda_+,\lambda_-}=\mathbf C\cdot 1\oplus \mathbf C\cdot\tau,
 \qquad
 \tau^2=\delta^2.
\end{equation}
The line \(\mathbf C\tau\) is the \emph{branch line}.
\end{construction}

\begin{proof}
Expanding the quadratic relation for \(K\) gives
\(K^2-(\lambda_++\lambda_-)K+\lambda_+\lambda_-=0\).  Hence
\[
(K-c)^2
=
K^2-2cK+c^2
=
\bigl((\lambda_++\lambda_-)K-\lambda_+\lambda_-\bigr)-2cK+c^2
=
c^2-\lambda_+\lambda_-
=
\delta^2.
\]
So every polynomial in \(K\) reduces uniquely to a scalar plus a
multiple of \(\tau\).
\end{proof}

\begin{lemma}[Universal branch moments; \ClaimStatusProvedHere]
\label{lem:universal-branch-moments}
In the rank-two branch algebra one has
\begin{equation}
\label{eq:universal-branch-moments}
K^m=a_m+b_m\tau,
\qquad
b_m=
\frac{\lambda_+^m-\lambda_-^m}{\lambda_+-\lambda_-},
\end{equation}
for every \(m\geq 1\).  Equivalently,
\begin{equation}
\label{eq:branch-line-projection}
\operatorname{pr}_{\tau}(K^m)
=
\frac{\lambda_+^m-\lambda_-^m}{\lambda_+-\lambda_-}\,\tau.
\end{equation}
\end{lemma}

\begin{proof}
Write \(K^m=a_m+b_m\tau\).  Evaluating on the two eigenspaces of
\(K\) yields
\[
\lambda_+^m=a_m+b_m\delta,
\qquad
\lambda_-^m=a_m-b_m\delta.
\]
Subtracting gives
\[
2\delta\,b_m=\lambda_+^m-\lambda_-^m.
\]
Since \(2\delta=\lambda_+-\lambda_-\), one obtains
\eqref{eq:universal-branch-moments}.
\end{proof}

\begin{definition}[Branch-line and primitive quotients]
\label{def:branch-line-primitive-quotients}
Let \(\cA\) be a modular Koszul chiral algebra whose spectral
characteristic factors through a rank-two sheet.
\begin{enumerate}[label=\textup{(\roman*)}]
\item The \emph{branch-line quotient} keeps only the scalar vertex data
  \(\kappa(\cA)\lambda_h\), the branch transport along edges, and the
  projection to the branch line \(\mathbf C\tau\).
\item The \emph{pure-branch primitive quotient} further discards every
  positive-genus vertex decoration and every separating edge.
\item The \emph{top-Hodge primitive quotient} keeps the top Hodge line
  \(\mathbf Q\lambda_{g(v)}\) on each vertex, retains only connected
  nonseparating graphs, and projects the spectral factor to the branch
  line.
\end{enumerate}
\end{definition}

\begin{definition}[Hodge depth]
\label{def:hodge-depth}
For a connected stable graph \(\Gamma\), define its
\emph{Hodge depth} by
\begin{equation}
\label{eq:hodge-depth}
\delta(\Gamma):=|E(\Gamma)|+\sum_{v\in V(\Gamma)}g(v).
\end{equation}
\end{definition}

\begin{theorem}[Depth formula; \ClaimStatusProvedHere]
\label{thm:hodge-depth-formula}
For every connected stable graph \(\Gamma\) of total genus \(g\),
\begin{equation}
\label{eq:depth-formula}
\delta(\Gamma)=g+|V(\Gamma)|-1.
\end{equation}
Consequently:
\begin{enumerate}[label=\textup{(\roman*)}]
\item \(\delta(\Gamma)\geq g\);
\item equality holds if and only if \(|V(\Gamma)|=1\).
\end{enumerate}
Thus the first primitive depth in genus \(g\) is exactly the one-vertex
sector.
\end{theorem}

\begin{proof}
The genus formula for a connected stable graph is
\[
 g=b_1(\Gamma)+\sum_{v}g(v)
 =|E(\Gamma)|-|V(\Gamma)|+1+\sum_v g(v).
\]
Rearranging gives \eqref{eq:depth-formula}.  The corollaries are
immediate.
\end{proof}

\section{Genus-\texorpdfstring{$2$}{2} transparency}
\label{sec:blr-genus-two-transparency}

The one-dimensional branch line forces a strong vanishing statement in
genus~\(2\).  The proof is entirely geometric: separating clutching sees
\(\lambda_2\), and nonseparating clutching kills it.

\begin{lemma}[Separating Hodge splitting; \ClaimStatusProvedHere]
\label{lem:separating-hodge-splitting}
For the separating clutching map
\[
\xi_{\mathrm{sep}}\colon
\overline{\mathcal M}_{1,1}\times\overline{\mathcal M}_{1,1}
\longrightarrow
\overline{\mathcal M}_{2},
\]
one has a canonical isomorphism of Hodge bundles
\[
\xi_{\mathrm{sep}}^*\mathbb E_2
\cong
\operatorname{pr}_1^*\mathbb E_1
\oplus
\operatorname{pr}_2^*\mathbb E_1.
\]
Hence
\begin{equation}
\label{eq:lambda2-separating-pullback}
\xi_{\mathrm{sep}}^*\lambda_2=\lambda_1\boxtimes\lambda_1.
\end{equation}
\end{lemma}

\begin{proof}
Let \(C=C_1\cup_q C_2\) be a separating nodal genus-\(2\) curve with
elliptic components \(C_1\) and \(C_2\).  The normalization description
of the dualizing sheaf identifies sections of \(\omega_C\) with pairs of
sections of \(\omega_{C_1}(q)\) and \(\omega_{C_2}(q)\) whose residues at
\(q\) add to zero.  On an elliptic curve every section of
\(\omega_{C_i}(q)\) is already holomorphic, so the residue condition is
automatic.  Therefore
\[
H^0(C,\omega_C)
\cong
H^0(C_1,\omega_{C_1})\oplus H^0(C_2,\omega_{C_2}),
\]
and this identification globalizes over the separating boundary family.
Taking top Chern classes gives
\eqref{eq:lambda2-separating-pullback}.
\end{proof}

\begin{lemma}[Nonseparating Hodge extension; \ClaimStatusProvedHere]
\label{lem:nonseparating-hodge-extension}
For the nonseparating clutching map
\[
\xi_{\mathrm{ns}}\colon
\overline{\mathcal M}_{1,2}
\longrightarrow
\overline{\mathcal M}_2,
\]
one has an exact sequence
\[
0\longrightarrow \mathbb E_1
\longrightarrow \xi_{\mathrm{ns}}^*\mathbb E_2
\longrightarrow \mathcal O
\longrightarrow 0.
\]
Consequently,
\begin{equation}
\label{eq:lambda2-nonseparating-pullback}
\xi_{\mathrm{ns}}^*\lambda_2=0.
\end{equation}
\end{lemma}

\begin{proof}
A point of the nonseparating boundary is an elliptic curve \(E\) with two
marked points \(p,q\) identified.  Sections of the dualizing sheaf of
the nodal curve correspond to meromorphic differentials on \(E\) with at
most simple poles at \(p\) and \(q\) and opposite residues.  Such a space
is generated by the holomorphic differential on \(E\) together with a
normalized differential of the third kind having residues \(+1\) and
\(-1\) at \(p\) and \(q\).  The quotient line is therefore trivial, which
gives the exact sequence above.  Since the quotient has vanishing first
Chern class, one gets
\[
\xi_{\mathrm{ns}}^*\lambda_2
=c_2(\xi_{\mathrm{ns}}^*\mathbb E_2)=0.
\]
\end{proof}

\begin{theorem}[Genus-\(2\) transparency on a one-line branch quotient; \ClaimStatusProvedHere]
\label{thm:genus-two-transparency}
Let \(R=S\oplus T\) be a commutative algebra such that \(T\) is
one-dimensional and \(T^2\subseteq S\).  Let
\(\Xi_g\in H^\bullet(\overline{\mathcal M}_g,\mathbf Q)\otimes R\) be a
clutching-compatible tower,
\[
\xi_{g_1,g_2}^*\Xi_{g_1+g_2}=\Xi_{g_1}\boxtimes\Xi_{g_2}.
\]
Assume that:
\begin{enumerate}[label=\textup{(\roman*)}]
\item the genus-\(1\) traceless part vanishes;
\item the genus-\(2\) traceless part lies on the top Hodge line:
  \begin{equation}
  \label{eq:genus-two-top-hodge-ansatz}
  \operatorname{pr}_T(\Xi_2)=\nu_2\,\tau\otimes\lambda_2
  \end{equation}
  for some nonzero generator \(\tau\) of \(T\).
\end{enumerate}
Then \(\nu_2=0\).
\end{theorem}

\begin{proof}
Restrict to the separating boundary.  By clutching,
\[
\xi_{\mathrm{sep}}^*\Xi_2=\Xi_1\boxtimes \Xi_1.
\]
Because the genus-\(1\) traceless part vanishes, the traceless part of
\(\Xi_1\boxtimes\Xi_1\) vanishes as well.  On the other hand,
\eqref{eq:genus-two-top-hodge-ansatz} and
Lemma~\ref{lem:separating-hodge-splitting} give
\[
\operatorname{pr}_T\bigl(\xi_{\mathrm{sep}}^*\Xi_2\bigr)
=
\nu_2\,\tau\otimes \lambda_1\boxtimes\lambda_1.
\]
Now \(\tau\neq 0\), and
\(\lambda_1\boxtimes\lambda_1\neq 0\) on
\(\overline{\mathcal M}_{1,1}\times\overline{\mathcal M}_{1,1}\), so
\(\nu_2=0\).  Lemma~\ref{lem:nonseparating-hodge-extension} shows that
there is no nonseparating contribution on the top Hodge line, because
\(\xi_{\mathrm{ns}}^*\lambda_2=0\).
\end{proof}

\begin{corollary}[Virasoro genus-\(2\) coefficient vanishes in the one-line quotient; \ClaimStatusProvedHere]
\label{cor:vir-genus-two-vanishing}
In every Virasoro branch-line quotient whose genus-\(2\) traceless part
lies on \(\lambda_2\), the coefficient \(\mu_2(c)\) vanishes
identically.
\end{corollary}

\begin{proof}
Apply Theorem~\ref{thm:genus-two-transparency} to the one-dimensional
branch line of the shared rank-two sheet.
\end{proof}

\begin{corollary}[The first primitive traceless coefficient begins in genus \texorpdfstring{$3$}{3}; \ClaimStatusProvedHere]
\label{cor:first-primitive-genus-three}
In the top-Hodge branch-line quotient, the first possible nonzero
primitive traceless contribution occurs in genus~\(3\).
\end{corollary}

\begin{proof}
Genus~\(1\) is scalar by hypothesis, and genus~\(2\) is annihilated by
Theorem~\ref{thm:genus-two-transparency}.
\end{proof}

\section{The pure-branch primitive genus-\texorpdfstring{$3$}{3} coefficient}
\label{sec:blr-pure-branch-g3}

The pure-branch quotient keeps only genus-zero vertices and
nonseparating edges.  In genus~\(3\) there is only one primitive graph:
the three-loop rose.

\begin{lemma}[Uniqueness of the primitive rose in genus \texorpdfstring{$3$}{3}; \ClaimStatusProvedHere]
\label{lem:genus-three-rose-unique}
In the pure-branch primitive quotient, the connected genus-\(3\)
sector is one-dimensional, generated by the one-vertex three-loop rose.
If
\[
\xi_{\mathrm{rose}}^{(3)}\colon
\overline{\mathcal M}_{0,6}
\longrightarrow
\overline{\mathcal M}_3
\]
glues the ordered pairs \((1,2)\), \((3,4)\), and \((5,6)\), then the
corresponding class is
\begin{equation}
\label{eq:rho-three-definition}
\rho_3:=\bigl(\xi_{\mathrm{rose}}^{(3)}\bigr)_*(1).
\end{equation}
\end{lemma}

\begin{proof}
In the pure-branch primitive quotient every vertex has genus~\(0\), so
the total genus is the first Betti number:
\[
3=b_1(\Gamma)=|E(\Gamma)|-|V(\Gamma)|+1.
\]
The primitive sector at minimal depth has \(|V(\Gamma)|=1\) by
Theorem~\ref{thm:hodge-depth-formula}, hence \(|E(\Gamma)|=3\).  The
unique connected one-vertex graph with three edges is the rose.
\end{proof}

\begin{theorem}[Pure-branch primitive coefficient on a rank-two sheet; \ClaimStatusProvedHere]
\label{thm:pure-branch-primitive-coefficient}
Let \(\cA\) be a modular Koszul chiral algebra whose spectral
characteristic factors through the rank-two sheet
\(R_{\lambda_+,\lambda_-}\).  In the pure-branch primitive quotient, the
first primitive traceless coefficient is
\begin{equation}
\label{eq:pure-branch-primitive-general}
\Theta^{\mathrm{pb}}_{3,\mathrm{prim}}(\cA)
=
\frac{\lambda_+^3-\lambda_-^3}{\lambda_+-\lambda_-}
\,\tau\otimes \rho_3.
\end{equation}
For the shared sheet
\(\Delta(x)=(1-3x)(1+x)\), this becomes
\begin{equation}
\label{eq:pure-branch-primitive-seven}
\Theta^{\mathrm{pb}}_{3,\mathrm{prim}}=7\,\tau\otimes \rho_3.
\end{equation}
\end{theorem}

\begin{proof}
By Lemma~\ref{lem:genus-three-rose-unique}, the primitive quotient is
one-dimensional on the graph side, so the only possible coefficient is a
scalar multiple of \(\tau\otimes \rho_3\).  Each loop contributes one
copy of the branch transport operator \(K\), so the total transport is
\(K^3\).  Projecting to the branch line and applying
Lemma~\ref{lem:universal-branch-moments} gives
\eqref{eq:pure-branch-primitive-general}.  On the shared sheet,
\((\lambda_+,\lambda_-)=(3,-1)\), hence the coefficient is
\((27-(-1))/4=7\).
\end{proof}

\begin{remark}[Rose coefficients in all genera]
\label{rem:rose-coefficients-all-genera}
The same proof gives the entire pure rose tower.  If
\(\rho_g^{\mathrm{rose}}\) denotes the class of the \(g\)-loop rose, then
\begin{equation}
\label{eq:rose-coefficients-all-genera}
\Theta^{\mathrm{pb}}_{g,\mathrm{rose}}(\cA)
=
\frac{\lambda_+^g-\lambda_-^g}{\lambda_+-\lambda_-}
\,\tau\otimes \rho_g^{\mathrm{rose}}.
\end{equation}
For the shared sheet this coefficient is
\(d_g=(3^g-(-1)^g)/4\).
\end{remark}

\section{The top-Hodge primitive quotient}
\label{sec:blr-top-hodge-primitive}

The pure-branch quotient isolates the universal constant\;\(7\), but it
forgets the scalar vertex data.  The next exact quotient restores the
top Hodge line on each vertex.

\begin{definition}[Graph-normalized primitive classes]
\label{def:graph-normalized-primitive-classes}
For \(0\leq h\leq g\), let
\[
\xi_h^{(g-h)}\colon
\overline{\mathcal M}_{h,2(g-h)}
\longrightarrow
\overline{\mathcal M}_g
\]
be the map that glues the ordered pairs
\((1,2),\dots,(2(g-h)-1,2(g-h))\).  Define
\begin{equation}
\label{eq:eta-gh-definition}
\bar\eta_{g,h}
:=
\frac{1}{2^{g-h}(g-h)!}
\bigl(\xi_h^{(g-h)}\bigr)_*(\lambda_h),
\qquad
\lambda_0:=1.
\end{equation}
The prefactor is the standard reciprocal automorphism factor for the
one-vertex \((g-h)\)-loop graph.
\end{definition}

\begin{remark}[The genus-\(3\) basis]
\label{rem:genus-three-basis-notation}
For genus~\(3\), the classes of
Definition~\ref{def:graph-normalized-primitive-classes} are
\[
\bar\rho_3:=\bar\eta_{3,0},
\qquad
\bar\epsilon_3:=\bar\eta_{3,1},
\qquad
\bar\gamma_3:=\bar\eta_{3,2}.
\]
These are respectively the three-loop rational rose, the elliptic
bi-rose with a \(\lambda_1\)-insertion, and the genus-\(2\) monorose
with a \(\lambda_2\)-insertion.
\end{remark}

\begin{theorem}[First primitive top-Hodge layer; \ClaimStatusProvedHere]
\label{thm:first-primitive-top-hodge-layer}
Let \(\cA\) be a modular Koszul chiral algebra whose spectral
characteristic factors through the rank-two sheet
\(R_{\lambda_+,\lambda_-}\).  In the top-Hodge primitive quotient, the
first primitive layer in genus \(g\geq 2\) is
\begin{equation}
\label{eq:primitive-top-hodge-general}
\Sigma_g(\cA)
=
\tau\otimes
\left(
\frac{\lambda_+^g-\lambda_-^g}{\lambda_+-\lambda_-}
\,\bar\eta_{g,0}
+
\kappa(\cA)
\sum_{h=1}^{g-1}
\frac{\lambda_+^{g-h}-\lambda_-^{g-h}}{\lambda_+-\lambda_-}
\,\bar\eta_{g,h}
\right).
\end{equation}
\end{theorem}

\begin{proof}
By Theorem~\ref{thm:hodge-depth-formula}, the first primitive depth in
genus \(g\) is the one-vertex sector.  A one-vertex primitive graph is
therefore determined by its vertex genus \(h\) and its loop number
\(g-h\).  On the graph side this contributes the class
\(\bar\eta_{g,h}\).  On the scalar side, a genus-\(h\) vertex contributes
\(\kappa(\cA)\lambda_h\) for \(h\geq 1\), whereas the genus-zero vertex
is the undeformed tree-level base and contributes only the unit.  On the
spectral side, each loop contributes one factor of \(K\), so the total
branch transport is \(K^{g-h}\).  Projecting to the branch line and
using Lemma~\ref{lem:universal-branch-moments} gives exactly
\eqref{eq:primitive-top-hodge-general}.
\end{proof}

\begin{corollary}[The genus-\(3\) primitive coefficient; \ClaimStatusProvedHere]
\label{cor:genus-three-primitive-top-hodge}
For genus~\(3\) one has
\begin{equation}
\label{eq:genus-three-primitive-top-hodge-general}
\Sigma_3(\cA)
=
\tau\otimes
\left(
\frac{\lambda_+^3-\lambda_-^3}{\lambda_+-\lambda_-}\,\bar\rho_3
+
\kappa(\cA)
\frac{\lambda_+^2-\lambda_-^2}{\lambda_+-\lambda_-}\,\bar\epsilon_3
+
\kappa(\cA)
\bar\gamma_3
\right).
\end{equation}
On the shared sheet \((\lambda_+,\lambda_-)=(3,-1)\), this becomes
\begin{equation}
\label{eq:genus-three-primitive-top-hodge-seven}
\Sigma_3(\cA)
=
\tau\otimes
\bigl(7\bar\rho_3+2\kappa(\cA)\bar\epsilon_3+\kappa(\cA)\bar\gamma_3\bigr).
\end{equation}
For Virasoro, where
\(\kappa(\mathrm{Vir}_c)=c/2\) by
Theorem~\ref{thm:modular-characteristic}, one gets
\begin{equation}
\label{eq:genus-three-primitive-virasoro}
\Sigma_3(\mathrm{Vir}_c)
=
\tau\otimes
\left(
7\bar\rho_3+c\bar\epsilon_3+\frac{c}{2}\bar\gamma_3
\right).
\end{equation}
\end{corollary}

\begin{proof}
Substitute \(g=3\) into
Theorem~\ref{thm:first-primitive-top-hodge-layer}.  For the shared
sheet,
\[
\frac{3^3-(-1)^3}{3-(-1)}=7,
\qquad
\frac{3^2-(-1)^2}{3-(-1)}=2,
\qquad
\frac{3-(-1)}{3-(-1)}=1.
\]
The Virasoro specialization follows from the scalar identity
\(\kappa(\mathrm{Vir}_c)=c/2\).
\end{proof}

\begin{remark}[What genus~\(3\) remembers]
\label{rem:what-genus-three-remembers}
Equation~\eqref{eq:genus-three-primitive-top-hodge-seven} separates the
first primitive traceless coefficient into three distinct sources:
\begin{enumerate}[label=\textup{(\roman*)}]
\item the universal pure-branch constant~\(7\), already visible on the
  three-loop rational rose;
\item the genus-\(1\) scalar shadow transported through two loops;
\item the genus-\(2\) scalar shadow transported through one loop.
\end{enumerate}
The quotient therefore exhibits a clean operator-level decomposition of
what the first primitive non-scalar correction can see.
\end{remark}

\section{Shared-sheet specializations and control cases}
\label{sec:blr-shared-sheet-specializations}

\begin{corollary}[Universal coefficients on the shared sheet; \ClaimStatusProvedHere]
\label{cor:shared-sheet-universal-coefficients}
On the shared sheet of
Theorem~\ref{thm:ds-bar-gf-discriminant}, the branch-line moments are
\begin{equation}
\label{eq:shared-sheet-dm}
 d_m:=\frac{3^m-(-1)^m}{4},
\end{equation}
so that
\[
 d_1=1,
 \qquad
 d_2=2,
 \qquad
 d_3=7,
 \qquad
 d_{m+2}=2d_{m+1}+3d_m.
\]
Accordingly,
\begin{equation}
\label{eq:shared-sheet-general-primitive-layer}
\Sigma_g(\cA)
=
\tau\otimes
\left(
 d_g\,\bar\eta_{g,0}
 +
 \kappa(\cA)
 \sum_{h=1}^{g-1}d_{g-h}\,\bar\eta_{g,h}
 \right)
\end{equation}
for every algebra on the shared sheet.
\end{corollary}

\begin{proof}
This is the specialization of
Lemma~\ref{lem:universal-branch-moments} and
Theorem~\ref{thm:first-primitive-top-hodge-layer} to
\((\lambda_+,\lambda_-)=(3,-1)\).  The recurrence follows from the
minimal polynomial \(K^2-2K-3=0\).
\end{proof}

\begin{example}[The scalar-only Heisenberg frame]
\label{ex:heisenberg-scalar-only-frame}
The Heisenberg frame example in Part~1 shows the opposite extreme from
the shared interacting sheet.  Its higher-genus package is already
controlled by the scalar invariant \(\kappa\), so there is no nontrivial
rank-two branch-line quotient to extract.  In that sense the Heisenberg
example is the scalar atom, whereas the shared
\(\widehat{\mathfrak{sl}}_2\)/Virasoro/\(\beta\gamma\) sheet is the
first genuine spectral atom.
\end{example}

\section{What remains open above the quotient}
\label{sec:blr-open-above-quotient}

The quotients constructed here are exact, but they are not the whole
story.  They record only what is forced by the scalar shadow, the
rank-two spectral sheet, and the stable-graph clutching algebra.  The
full modular deformation complex contains strictly more information.

\begin{problem}[Branch realization of the full cyclic complex]
\label{prob:branch-realization-full-cyclic}
Construct a natural morphism from the full cyclic deformation complex
\[
\mathcal R_{\mathrm{br}}\colon
\Defcyc(\cA)
\longrightarrow
\operatorname{End}(V_{\mathrm{br}})\widehat\otimes
C^*(\text{stable graphs})
\]
whose logarithmic connected trace recovers the spectral potential
\(\Psi_{\cA}(x)\).
\end{problem}

\begin{conjecture}[Spectral shadow principle; \ClaimStatusConjectured]
\label{conj:spectral-shadow-principle}
For every modular Koszul chiral algebra~$\cA$, the full universal
class~$\Theta_{\cA}$ admits a natural branch realization
\[
\mathcal R_{\mathrm{br}}\colon
L_{\cA}^\circ
\longrightarrow
\operatorname{End}(V_{\mathrm{br},\cA})\llbracket x\rrbracket
\]
whose connected one-loop image is exactly
\[
-\log\Delta_{\cA}(x) = \Psi_{\cA}(x).
\]
The proved spectral potential is the one-loop
connected shadow of the full graph-valued Maurer--Cartan element.
\end{conjecture}

\begin{remark}[Why the logarithm is canonical]
\label{rem:logarithm-canonical}
The conjecture is structurally forced: a Maurer--Cartan element in a
completed dg Lie algebra decomposes into connected graph contributions by
the standard exponential/logarithm correspondence.  The determinant
$\Delta_{\cA}$ is the partition function (disconnected sum), and
$\Psi_{\cA}=-\log\Delta_{\cA}$ is the connected potential.  The
identification would say that the spectral connected potential is exact
at the one-loop level, with all higher-loop connected contributions lying
in the traceless sector.
\end{remark}

\begin{problem}[Renormalization of primitive quotient coefficients]
\label{prob:renormalization-primitive-coefficients}
Determine whether the coefficients in the primitive quotients are exact
images of the full Maurer--Cartan class or receive corrections from
non-top-Hodge vertex decorations and genuinely higher cyclic
\(L_\infty\)-operations.  In particular, decide whether the quotient
formula
\[
\Sigma_3(\mathrm{Vir}_c)
=
\tau\otimes
\left(
7\bar\rho_3+c\bar\epsilon_3+\frac{c}{2}\bar\gamma_3
\right)
\]
is already the exact genus-\(3\) traceless image of
\(\Theta_{\mathrm{Vir}_c}\) under a canonical branch realization.
\end{problem}

\begin{problem}[Beyond top Hodge]
\label{prob:beyond-top-hodge}
Describe the first quotient above the top-Hodge primitive layer.  A
natural candidate is the quotient that keeps one additional tautological
vertex decoration, for example a \(\psi\)-class or a boundary divisor
class, together with the same branch-line projection.  Theorem~\ref{thm:hodge-depth-formula}
suggests that the first new terms should appear either at the same depth
with richer vertex decorations or one depth higher with multi-vertex
primitive graphs.
\end{problem}

\begin{remark}[The structural picture]
\label{rem:blr-structural-picture}
The quotient theory established here has a clean shape.
\begin{enumerate}[label=\textup{(\roman*)}]
\item The scalar Maurer--Cartan skeleton is the abelian shadow of the
  full universal class.
\item The spectral discriminant produces a connected cumulant tower.
\item The first nontrivial correction to a scalar shadow is forced into a
  traceless quadratic sector.
\item In the Virasoro family, the positive-weight sector is contractible,
  so the first finite obstruction is spectral.
\item On every one-line branch quotient with scalar square, genus~\(2\)
  is transparent.
\item The first primitive top-Hodge coefficient therefore occurs in
  genus~\(3\), and on the shared sheet its universal pure-branch part is
  the constant~\(7\).
\end{enumerate}
This is the exact content of the branch-line and primitive quotient
programme extracted from the thread.
\end{remark}
